%! The Normal Distribution and z-Scores — Unit 1, Lesson 1.7 — Slides (Beamer)
\documentclass[aspectratio=169, 11pt]{beamer}
\usepackage{apstats-beamer}   % provides \navyheader{} and \sectionlabel[color]{}
\usetikzlibrary{positioning}  % -beamer loads tikz but not this library

% The C and Y column types live in apstats-article, which a beamer deck does not
% load — redeclare them here rather than pulling in the article preamble.
\newcolumntype{C}[1]{>{\centering\arraybackslash}p{#1}}
\newcolumntype{Y}{>{\centering\arraybackslash}X}

\begin{document}

% TITLE SLIDE (CourseName is not defined in beamer, so the name is literal)
{
\setbeamercolor{background canvas}{bg=navy}
\begin{frame}[plain]
  \vspace{2.8cm}\hspace{0.55cm}%
  \begin{minipage}{0.9\textwidth}
    {\color{goldacc}\bfseries\small LESSON 1.7}\\[0.5cm]
    {\color{white}\bfseries\LARGE The Normal Distribution and z-Scores}\\[1.2cm]
    {\color{skymid}\small Statistics~$\cdot$~Unit 1}
  \end{minipage}
\end{frame}
}

% GRADUAL RELEASE deck: title → targets → warm-up hook → I DO divider → ONE FRAME PER NOTES ROW
% (parameter/statistic · z-score · normal + empirical rule · the crux) → WE DO (Guided Practice,
% The Bread Loaf) → spoken debrief → close & START THE HOMEWORK, which is the individual
% practice — there is no independent practice set in the notes → AP Practice → close.
% There is NO group activity in this course and NO spoiler rule: vocabulary is given up front,
% and the notes frames ARE the direct instruction. This is the LAST content lesson of Unit 1,
% so the closing frame looks ahead to the unit test and to Unit 2.
\newcommand{\redink}[1]{{\color{redacc}\bfseries #1}}

% ── LEARNING TARGETS ─────────────────────────────────────────────────────────
\begin{frame}[t, plain]
  \navyheader{Today's targets}
  \sectionlabel{Lesson 1.7}
  By the end of today, I can\ldots
  \begin{itemize}\setlength{\itemsep}{5pt}
    \item tell a \textbf{parameter} from a \textbf{statistic}, and write the right symbol.
    \item compute a \textbf{$z$-score} and interpret it in context --- sign included.
    \item use the \textbf{empirical rule} on an approximately \textbf{normal} distribution.
    \item use $z$-scores to compare two results from two different distributions.
  \end{itemize}
  \begin{block}{How today works}
    Warm-up $\to$ \textbf{notes together} $\to$ \textbf{one we work as a class} $\to$
    debrief $\to$ \textbf{you start the homework here, alone}.
  \end{block}
\end{frame}

% ── HOOK (off the warm-up) ───────────────────────────────────────────────────
{
\setbeamercolor{background canvas}{bg=navy}
\begin{frame}[plain]
  \vspace{1.4cm}
  \begin{center}
    {\color{skymid}\bfseries\small ONE SCHOLARSHIP. TWO ATHLETES.}\\[0.9cm]
    {\color{white}\bfseries\Large Maya threw \textbf{12.0}. \quad Dev ran \textbf{12.3}.\\[0.4cm]
    Dev's number is bigger. Dev did not win.}\\[0.9cm]
    {\color{goldacc}\small Metres and seconds. Two scales that will not sit on the same page ---
    until you count in the right units. Warm-up 2 already did the arithmetic.}
  \end{center}
\end{frame}
}

% ── I DO — direct instruction begins ─────────────────────────────────────────
{
\setbeamercolor{background canvas}{bg=navy}
\begin{frame}[plain]
  \vspace{1.8cm}\hspace{0.8cm}%
  \begin{minipage}{0.86\textwidth}
    {\color{goldacc}\bfseries\small GUIDED NOTES \ \textbullet\ \ I DO}\\[0.7cm]
    {\color{white}\bfseries\Large Open to your Guided Notes page. Fill each blank as we
    name it --- not before.}\\[0.9cm]
    {\color{skymid}\small Five terms go in the vocabulary box today. Write each one the
    moment we use it.}
  \end{minipage}
\end{frame}
}

% ── NOTES ROW 1 ──────────────────────────────────────────────────────────────
\begin{frame}[t, plain]
  \navyheader{Parameter or statistic?}
  \sectionlabel{notes, Parameter \& Statistic $\cdot$ everybody, or just some of them}
  \vspace{4pt}
  \begin{columns}[t]
    \begin{column}{0.47\textwidth}
      \begin{block}{Parameter}
        A numerical summary of a \textbf{population} --- everybody (EK VAR-2.A.1).
      \end{block}
      \vspace{4pt}
      {\small Population mean $\mu$ \quad$\cdot$\quad population standard deviation $\sigma$
      \\[3pt]
      \emph{All 500 loaves the bakery weighed.}}
    \end{column}
    \begin{column}{0.47\textwidth}
      \begin{block}{Statistic}
        A numerical summary of a \textbf{sample} --- some of them (Lesson 1.5).
      \end{block}
      \vspace{4pt}
      {\small Sample mean $\bar x$ \quad$\cdot$\quad sample standard deviation $s_x$\\[3pt]
      \emph{The 40 loaves the inspector weighed.}}
    \end{column}
  \end{columns}
  \vspace{12pt}
  \begin{center}
    {\large Ask \redink{who got measured} --- everyone, or a handful?}\\[6pt]
    {\small\color{slate} Same idea, different object. The symbol you write depends on the answer.
    Problems 1--3.}
  \end{center}
\end{frame}

% ── NOTES ROW 2 ──────────────────────────────────────────────────────────────
\begin{frame}[t, plain]
  \navyheader{The z-score}
  \sectionlabel{notes, $z$-score $\cdot$ stop counting points, start counting standard deviations}
  \vspace{6pt}
  \begin{center}
    {\Large $z = \dfrac{x_i - \mu}{\sigma}$}
  \end{center}
  \vspace{8pt}
  \begin{columns}[t]
    \begin{column}{0.47\textwidth}
      {\small
      % \redink is text-mode (\bfseries) — inside math, colour with \mathbf instead.
      \textbf{Maya:} $\dfrac{12.0 - 9.0}{1.5} = {\color{redacc}\mathbf{2.0}}$\\[8pt]
      \textbf{Dev:} $\dfrac{12.3 - 13.2}{0.6} = {\color{redacc}\mathbf{-1.5}}$}
    \end{column}
    \begin{column}{0.47\textwidth}
      \begin{itemize}\setlength{\itemsep}{5pt}
        \item Subtract the mean, then divide by that distribution's \textbf{own} standard
              deviation --- metres over metres, seconds over seconds, so the \redink{units cancel}
              (EK VAR-2.B.1, 2.B.2).
        \item $z < 0$: below the mean. $z = 0$: exactly at it.
        \item Say it in context: \emph{``1.5 standard deviations below the mean 100 m time''} ---
              which in a race is \textbf{good}.
      \end{itemize}
    \end{column}
  \end{columns}
  \vspace{6pt}
  \begin{center}
    {\small\color{slate} Problems 4--6. I work both $z$-scores at the board; you work the rest.
    Problem 5 wants the direction \emph{and} what it means in a sprint.}
  \end{center}
\end{frame}

% ── NOTES ROW 3 ──────────────────────────────────────────────────────────────
\begin{frame}[t, plain]
  \navyheader{The normal distribution and the empirical rule}
  \sectionlabel{notes, Empirical Rule $\cdot$ mound-shaped, symmetric, fixed by two numbers}
  \vspace{2pt}
  \begin{columns}[t]
    \begin{column}{0.52\textwidth}
      \begin{center}
      \begin{tikzpicture}[scale=0.72]
        \draw[navy, very thick] plot[smooth] coordinates
          {(-3.5,0.01) (-3,0.03) (-2.5,0.11) (-2,0.34) (-1.5,0.81) (-1,1.52) (-0.5,2.21)
           (0,2.50) (0.5,2.21) (1,1.52) (1.5,0.81) (2,0.34) (2.5,0.11) (3,0.03) (3.5,0.01)};
        \foreach \x/\h in {-3/0.03, -2/0.34, -1/1.52, 0/2.50, 1/1.52, 2/0.34, 3/0.03} {
          \draw[linegray, thin, dashed] (\x,0) -- (\x,\h);}
        \draw[->, thick] (-4.0,0) -- (4.0,0);
        \foreach \x/\lab in {-3/$-3$, -2/$-2$, -1/$-1$, 0/$\mu$, 1/$+1$, 2/$+2$, 3/$+3$} {
          \draw (\x,-0.08) -- (\x,0.08); \node[below, font=\scriptsize] at (\x,-0.1) {\lab};}
        \node[font=\scriptsize] at (0,-0.68) {standard deviations from the mean};
      \end{tikzpicture}
      \end{center}
    \end{column}
    \begin{column}{0.44\textwidth}
      \begin{itemize}\setlength{\itemsep}{6pt}
        \item \textbf{Mound-shaped and symmetric}, and completely fixed by $\mu$ and $\sigma$
              (EK VAR-2.A.2).
        \item Real data is never exactly this. We say \redink{approximately} normal
              (EK VAR-2.A.4).
        \item \textbf{The empirical rule} (EK VAR-2.A.3): about \redink{68\%} within 1 standard
              deviation, \redink{95\%} within 2, \redink{99.7\%} within 3.
      \end{itemize}
    \end{column}
  \end{columns}
  \vspace{6pt}
  \begin{center}
    {\small\color{slate} Problem 7 first --- 500 loaves: 68.2, 95.2, 99.8. \quad 400
    thermometers: 68.25, 95.25, 99.75. \quad Same shape, same numbers.}
  \end{center}
\end{frame}

% ── NOTES ROW 4 — THE CRUX ───────────────────────────────────────────────────
\begin{frame}[t, plain]
  \navyheader{Why 12.3 lost to 12.0}
  \sectionlabel{notes, Whose Day Was Better? $\cdot$ the whole point of today}
  \vspace{6pt}
  \begin{itemize}\setlength{\itemsep}{8pt}
    \item A raw number means nothing on its own. It only means something \redink{next to the
          distribution it came from} (EK VAR-2.C.1).
    \item Subtracting the mean was not enough --- $+3.0$ metres and $-0.9$ seconds are still two
          different units. Dividing by each field's own standard deviation is the fix.
    \item \textbf{And keep the sign.} $-1.5$ does not mean ``worse.'' It means below the mean
          \emph{time} --- which in a race is \redink{faster than the field}.
  \end{itemize}
  \vspace{8pt}
  \begin{block}{But the rule is a fact about a shape}
    A skewed or two-peaked record does \textbf{not} hand you 68--95--99.7. Problem 14 is where
    you say why.
  \end{block}
\end{frame}

% ── WE DO ────────────────────────────────────────────────────────────────────
\begin{frame}[t, plain]
  \navyheader{Guided Practice: The Bread Loaf}
  \sectionlabel[goldacc]{We do $\cdot$ you hold the pen}
  Back to the bakery as a model: loaf weights are approximately normal, $\mu = 800$ g,
  $\sigma = 20$ g.
  \begin{itemize}\setlength{\itemsep}{5pt}
    \item A loaf weighs 834 g. Find its $z$-score --- then say what it \emph{means}, in context.
    \item Roughly what percent of loaves weigh more than 840 g? (Look at where 840 lands first.)
    \item A customer calls a 754 g loaf ``way underweight.'' Is the complaint fair?
  \end{itemize}
  \begin{block}{The question behind the question}
    Problems 19 and 20 ask you to \redink{judge}, not just compute --- and to name one honest
    limit of the model you judged with.
  \end{block}
\end{frame}

% ── DEBRIEF (spoken) ─────────────────────────────────────────────────────────
\begin{frame}[t, plain]
  \navyheader{Debrief: say it out loud}
  \sectionlabel{8 minutes $\cdot$ pens down, papers up --- nothing to write here}
  \vspace{6pt}
  \begin{columns}[t]
    \begin{column}{0.46\textwidth}
      \begin{block}{Both z-scores, and the direction}
        Computing $2.0$ and $-1.5$ is half an answer. Naming who had the better day, and why, is
        the other half.
      \end{block}
    \end{column}
    \begin{column}{0.46\textwidth}
      \begin{block}{Shape, not units}
        Grams and degrees, 500 and 400 --- and the same three percentages. It was never about
        what was being measured.
      \end{block}
    \end{column}
  \end{columns}
  \vspace{12pt}
  \begin{block}{Two things to say out loud}
    Interpret one $z$-score in full --- how many standard deviations, which side, of what.
    Then name a variable where a \redink{negative} $z$-score is good news, and one where it
    is bad.
  \end{block}
\end{frame}

% ── YOU DO — the homework is the individual practice, started in class ───────
\begin{frame}[t, plain]
  \navyheader{You do: start the homework}
  \sectionlabel{Last 8 minutes $\cdot$ on your own $\cdot$ finish it at home}
  \begin{itemize}\setlength{\itemsep}{5pt}
    \item Four summaries: parameter or statistic?
    \item Three heart rates, three $z$-scores --- and what the negative one means.
    \item Ana scored 82, Ben scored 88. Who did better?
    \item A bus route's travel times, and one tail to find.
    \item One multiple-choice item, with a one-sentence justification.
  \end{itemize}
  \begin{block}{Silent and alone --- I am circulating}
    \emph{Item 3} is the one I am reading over your shoulder. This is your \redink{scored}
    homework, on paper, \textbf{due the first class after two study halls}. Start with 1, 2,
    and 3 while I am still here to answer questions.
  \end{block}
\end{frame}

% ── AP PRACTICE (assigned, unscored) ────────────────────────────────────────
\begin{frame}[t, plain]
  \navyheader{AP Practice --- assigned, not scored}
  \sectionlabel[goldacc]{Newborn birth weights $\cdot$ bring it to the unit-test review}
  \begin{itemize}\setlength{\itemsep}{5pt}
    \item Four multiple-choice items on newborn birth weights --- including one where the right
          answer is that the empirical rule \textbf{does not apply}.
    \item One free-response set: compute, interpret, compare across two hospitals, and judge
          whether the model fits at all.
    \item Optional: the SAT and the ACT, put on the same scale.
  \end{itemize}
  \begin{block}{Not graded --- but this is what the exam looks like}
    The AP Practice page is \textbf{not scored}. Work it for the reps and bring it to the unit-test
    review. The \redink{graded} homework is the section behind it. And on the exam: in context,
    or it does not count.
  \end{block}
\end{frame}

% ── CLOSE — last content lesson of Unit 1 ───────────────────────────────────
{
\setbeamercolor{background canvas}{bg=navy}
\begin{frame}[plain]
  \vspace{1.2cm}\hspace{0.8cm}%
  \begin{minipage}{0.86\textwidth}
    {\color{goldacc}\bfseries\small BEFORE YOU GO}\\[0.6cm]
    {\color{white}\bfseries\Large Points, seconds, grams, and degrees cannot be compared.
    Standard deviations can. Once you count in standard deviations, every distribution speaks the
    same language --- and if it is mound-shaped, it always says 68, 95, and 99.7.}\\[0.8cm]
    {\color{skymid}\small That is the last of Unit 1. Next: the \textbf{Unit 1 test} --- variables
    and their displays, shape, centre, spread, resistance, boxplots, comparing two groups, and
    today's normal model. Then Unit 2 --- Exploring Two-Variable Data: one variable at a time is
    over, and we ask whether two of them move together.}
  \end{minipage}
\end{frame}
}

\end{document}
